\subsection{Amostragem estratificada e estimação}
Considere uma população finita $U$, de tamanho $N$, particionada em $H$ estratos disjuntos $U_h$, com tamanhos $N_h$. Em cada estrato, uma amostra aleatória simples sem reposição $s_h$ contém $n_h$ unidades. A estratificação permite controlar a composição do sorteio por características conhecidas da população; não exige frações de seleção iguais entre estratos. Esses fundamentos são tratados por \citeonline{cochran1977} e \citeonline{lohr2022}.

Para a unidade $i$ pertencente ao estrato $h$, a probabilidade de inclusão e o peso de desenho são
\begin{equation}\label{eq:inclusao}
 \pi_i=\frac{n_h}{N_h},\qquad w_i=\frac{1}{\pi_i}=\frac{N_h}{n_h}.
\end{equation}
Um estrato incluído integralmente tem $n_h=N_h$ e peso um. Para uma variável $z_i$, a média populacional pode ser estimada por
\begin{equation}\label{eq:mediaestratificada}
 \widehat{\overline Z}_{\mathrm{est}}
 =\sum_{h=1}^{H}\frac{N_h}{N}\overline z_h
 =\frac{1}{N}\sum_{i\in s}w_i z_i,
 \qquad\overline z_h=\frac{1}{n_h}\sum_{i\in s_h}z_i.
\end{equation}
Nessa alocação fixa, $\sum_{i\in s}w_i=N$. Sob sorteios independentes entre estratos, a variância estimada da média é
\begin{equation}\label{eq:varianciaestratificada}
 \widehat{\operatorname{Var}}_d(\widehat{\overline Z}_{\mathrm{est}})
 =\sum_{h=1}^{H}\left(\frac{N_h}{N}\right)^2
 \left(1-\frac{n_h}{N_h}\right)\frac{s_{z,h}^2}{n_h},
\end{equation}
em que $s_{z,h}^2$ é a variância amostral no estrato; estratos integrais contribuem com zero. A correção para população finita é relevante quando a fração amostral não é desprezível \cite{cochran1977}. Essas expressões descrevem o sorteio inicial, não o mecanismo posterior de recuperação de letras.

Se $A_i$ indica letra aceita, a média ponderada entre aceitas assume a forma de razão
\begin{equation}\label{eq:mediaaceitas}
 \widehat{\overline Z}_A=\frac{\sum_{i\in s}w_i A_i z_i}{\sum_{i\in s}w_i A_i}.
\end{equation}
Ela descreve a composição do domínio aceito segundo os pesos originais, mas não estima automaticamente a média de $z$ em todo $U$. O evento $A_i=1$ não é uma segunda seleção probabilística com probabilidades conhecidas. Disponibilidade, ambiguidade e qualidade de correspondência devem ser distinguidas de inclusão amostral, em consonância com a discussão de cobertura e não resposta de \citeonline{lohr2022}.

Para a razão da Equação~\ref{eq:mediaaceitas}, quando a variável está observada nos casos considerados, a linearização usa $u_i=A_i(z_i-\widehat{\overline Z}_A)$, incluindo zeros para as não aceitas. Se $\widehat N_A=\sum_{i\in s}w_iA_i$, a variância aproximada é
\begin{equation}\label{eq:varianciarazao}
 \widehat{\operatorname{Var}}_d(\widehat{\overline Z}_A)=\sum_h v_h,
 \qquad v_h=\frac{N_h^2(1-n_h/N_h)s_{u,h}^2}{n_h\widehat N_A^2}.
\end{equation}
Essa é uma aproximação de razão sob o desenho, distinta de tratar as aceitas como um novo sorteio aleatório simples \cite{statcan2010}. Nas comparações marginais, a média do catálogo é fixa, e a implementação utiliza graus de liberdade de Satterthwaite,
\begin{equation}\label{eq:dfrazao}
 \nu=\frac{(\sum_hv_h)^2}{\sum_{h:n_h<N_h}v_h^2/(n_h-1)},
\end{equation}
quando a variância é estimável e positiva. Células sem suporte para essa estimação não recebem conclusão automática de equivalência.

Compatibilidade marginal não se demonstra por deixar de rejeitar uma igualdade. Para uma diferença populacional $\theta$ e margem substantiva $\varepsilon>0$, a formulação de equivalência é
\begin{equation}\label{eq:equivalencia}
 H_0:\theta\le-\varepsilon\ \text{ou}\ \theta\ge\varepsilon,
 \qquad H_1:-\varepsilon<\theta<\varepsilon.
\end{equation}
No procedimento de dois testes unilaterais (TOST) com aproximação por $t$ de Student e nível $\eta$ em cada teste, a equivalência exige que
\begin{equation}\label{eq:tost}
 [\widehat\theta-t_{1-\eta,\nu}\widehat{\operatorname{SE}},
  \widehat\theta+t_{1-\eta,\nu}\widehat{\operatorname{SE}}]
 \subset(-\varepsilon,\varepsilon).
\end{equation}
Aqui $t_{q,\nu}$ é o quantil $q$ com $\nu$ graus de liberdade. Para $\eta=0,05$, usa-se intervalo marginal de 90\%. A lógica de equivalência e a necessidade de justificar margens são discutidas por \citeonline{wellek2010}. Equivalência, diferença apoiada e inconclusão são resultados distintos; concordância em algumas marginais não certifica representatividade multivariada ou identidade correta de letras.

\subsection{Regressão linear e regularização Ridge}
Seja $y_i$ a popularidade observada e $\bm x_i$ o vetor de $p$ características de uma faixa. A regressão linear aproxima a resposta por $\beta_0+\bm x_i^\top\bm\beta$. Com várias características correlacionadas, a estimação por mínimos quadrados pode produzir coeficientes instáveis. A regressão Ridge acrescenta uma penalização quadrática, introduzindo uma troca entre viés e variância \cite{hastie2009,james2023}.

Para o conjunto de treinamento $T$, os coeficientes são estimados por
\begin{equation}\label{eq:ridge}
 (\widehat\beta_0,\widehat{\bm\beta})=
 \underset{\beta_0,\bm\beta}{\operatorname{argmin}}
 \left\{\sum_{i\in T}(y_i-\beta_0-\bm x_i^\top\bm\beta)^2
 +\alpha\sum_{j=1}^{p}\beta_j^2\right\},\qquad\alpha>0.
\end{equation}
O intercepto não é penalizado. O parâmetro $\alpha$ controla o encolhimento dos coeficientes; valores maiores impõem penalização mais intensa. A proposta clássica de Ridge é de \citeonline{hoerl1970}. A comparação aqui mantém essa família nos quatro conjuntos de variáveis para examinar a informação incremental das representações.

Como a penalização depende da escala, cada variável numérica é padronizada no respectivo treino:
\begin{equation}\label{eq:padronizacao}
 x_{ij}^{*}=\frac{x_{ij}-\overline x_{j,T}}{s_{j,T}},\qquad
 s_{j,T}^2=\frac{1}{|T|}\sum_{i\in T}(x_{ij}-\overline x_{j,T})^2.
\end{equation}
A implementação usa desvio com divisor $|T|$. Uma coluna constante é centrada para zero. Médias, desvios e demais transformações são estimados apenas em $T$ e aplicados à avaliação, evitando que informação do teste participe do ajuste \cite{james2023}. Ridge otimiza quadrados penalizados; a escolha de sua penalização pode, contudo, utilizar outra métrica de validação, como o MAE.

\subsection{Medidas de erro e contribuição incremental}
Para $n$ observações de avaliação, denote por $\widehat y_i^{(m)}$ a previsão da configuração $m$, produzida sem usar o componente de $i$ no treino. Defina o erro de previsão $e_i^{(m)}=y_i-\widehat y_i^{(m)}$. O erro médio absoluto (MAE, \textit{mean absolute error}) e a raiz do erro quadrático médio (RMSE, \textit{root mean squared error}) são
\begin{align}
 \operatorname{MAE}_m&=\frac{1}{n}\sum_{i=1}^{n}|e_i^{(m)}|,\label{eq:mae}\\
 \operatorname{RMSE}_m&=\sqrt{\frac{1}{n}\sum_{i=1}^{n}(e_i^{(m)})^2}.\label{eq:rmse}
\end{align}
Essas definições são apresentadas por \citeonline{hyndman2021}. Ambas têm a unidade da resposta; o RMSE é mais sensível a desvios grandes devido ao quadrado. O MAE foi adotado como métrica primária de seleção e comparação.

Também é calculado o coeficiente de determinação nas previsões externas reunidas:
\begin{equation}\label{eq:r2}
 R_m^2=1-\frac{\sum_{i=1}^{n}(y_i-\widehat y_i^{(m)})^2}
 {\sum_{i=1}^{n}(y_i-\overline y)^2},\qquad
 \overline y=\frac{1}{n}\sum_{i=1}^{n}y_i.
\end{equation}
A definição compara a soma de erros quadráticos com a variação em torno da média \cite{james2023}. Fora do treino, $R^2$ pode ser negativo. Aqui ele resume os casos reunidos, não a média dos coeficientes dos cinco folds; não é uma proporção de variação causada pelas letras.

A diferença pareada e a redução relativa de MAE entre configurações $a$ e $b$ são
\begin{align}
 \Delta_{a-b}&=\frac{1}{n}\sum_{i=1}^{n}
 (|e_i^{(a)}|-|e_i^{(b)}|)=\operatorname{MAE}_a-\operatorname{MAE}_b,\label{eq:delta}\\
 G_{a:b}&=100\frac{\operatorname{MAE}_b-\operatorname{MAE}_a}{\operatorname{MAE}_b}.\label{eq:ganho}
\end{align}
Valor negativo de $\Delta$ favorece $a$; valor positivo de $G$ representa redução percentual do erro. O pareamento exige as mesmas faixas e previsões externas correspondentes. Uma sensibilidade adicional mantém esses ajustes e aplica os pesos originais apenas à avaliação:
\begin{equation}\label{eq:maeponderado}
 \operatorname{MAE}_{m,w}=\frac{\sum_i w_i|e_i^{(m)}|}{\sum_i w_i}.
\end{equation}
Essa escolha não equivale a treinar com pesos e não corrige automaticamente seleção por disponibilidade de letras.

\subsection{Descritores e embeddings de letras}
Descritores explícitos convertem propriedades observáveis de uma letra em números: extensão, diversidade de segmentos e repetição de linhas, por exemplo. Embeddings são vetores densos produzidos por um modelo pré-treinado. Os primeiros têm definições diretamente interpretáveis; os segundos oferecem uma representação aprendida de relações textuais. Nenhuma das duas abordagens garante adequação uniforme entre línguas ou ganho para a resposta de interesse.

\citeonline{parada2024lyrics} analisam características de letras em inglês ao longo do tempo. Seus descritores e sua base derivada do LFM2b \cite{parada2024data} motivaram a trajetória do projeto, mas não são as 21 medidas próprias aqui utilizadas. A distinção evita apresentar uma implementação diferente como reprodução daquele estudo.

O multilingual E5 foi proposto para embeddings multilingues \cite{wang2024}. Neste trabalho, o encoder permanece congelado e não recebe popularidade, artista ou título como entrada adicional. O tokenizer de subpalavras é distinto da segmentação lexical dos descritores. Métodos como SentencePiece permitem tokenização sem exigir segmentação prévia em palavras \cite{kudo2018}; isso não torna comparáveis, por definição, contagens lexicais de diferentes sistemas de escrita. A receita efetivamente aplicada ao corpus é especificada na Seção~\ref{subsec:representacoes}.

\subsection{Validação cruzada aninhada e dependência}
A validação cruzada estima desempenho fora do ajuste. Se a mesma avaliação escolhe hiperparâmetros e anuncia o resultado, a comparação pode ser otimista. No desenho aninhado, as partições internas selecionam hiperparâmetros e as externas avaliam o procedimento escolhido \cite{cawley2010}. A escolha de partições e a estimação de erro de generalização são discutidas em \citeonline{hastie2009}. Em português, \citeonline{izbicki2020} apresentam o aprendizado de máquina sob a perspectiva de risco preditivo, distinguindo avaliação de generalização de simples qualidade de ajuste.

Faixas do mesmo artista, versões relacionadas e letras idênticas podem compartilhar informação. Por isso, a unidade de separação deve reunir as relações identificadas, em vez de sortear cada linha de forma isolada. Os grupos adotados são componentes operacionais; ausência de sobreposição dessas chaves não prova independência estatística universal.

Há dependência adicional entre resultados de CV, pois os conjuntos de treino se sobrepõem. \citeonline{bengio2004} mostram os problemas de estimar a variância ignorando tais correlações. Consequentemente, os cinco folds não são tratados como cinco réplicas independentes, nem usados em um teste $t$ ingênuo. A análise auxiliar dos contrastes condiciona em ajustes já produzidos e tem alcance diferente da incerteza integral de todo o treinamento, conforme explicitado no Apêndice~\ref{ap:incerteza}.
